% ORIG03-SEGMENT-CODE: OR03-S001
% ORIG03-SEGMENT-ID: d90.orig.v1.tr03.seg0001
% ORIG03-MODULE-ID: map.assessment.0001
% SPDX-License-Identifier: CC-BY-SA-4.0
% Karya asli berbahasa Indonesia; tidak menyalin butir dari sumber donor.

\section{Peta asesmen dan kontrak topologi}
\label{orig03:section:peta-asesmen}

Bagian ini memetakan asesmen Original-03 tanpa mengulang bunyi soal. Seluruh
teks baru pada tranche ini ditulis secara independen dan dilisensikan di bawah
CC BY-SA 4.0. Notasi mengikuti spine Habring serta suplemen dan tranche yang
sudah diterima, tetapi tidak ada butir donor yang disalin. Materi program linear,
integer, simpleks, pemodelan riset operasi, dan wilayah O018 lainnya sengaja
tidak dimasukkan.

Kontrak yang dapat diurai mesin adalah sebagai berikut. Setiap butir induk
memiliki ID \texttt{*.problem.NNNN}. Setiap subbutir memiliki ID
\texttt{*.prompt.NNNN} dan tepat satu pasangan
\texttt{*.hint1.NNNN}/\texttt{*.hint2.NNNN}, satu
\texttt{*.answer.NNNN}, serta satu \texttt{*.solution.NNNN}. Komentar
\texttt{ORIG03-ASSESSMENT-MAP-ID} pada peta ini hanya merupakan rujukan
nondeklaratif. Deklarasi otoritatif \texttt{ORIG03-STABLE-ID} dan label LaTeX
yang berpadanan berada tepat pada butir soal di modul 01--10. Tujuh label
\texttt{orig03:rubric:proof:0001} sampai
\texttt{orig03:rubric:proof:0007} disediakan sebagai rujukan logis bagi rubrik
bukti yang dihimpun oleh agregator.

\subsection{Lima puluh empat prompt baru}

\paragraph{Diagnostik prasyarat: 20 prompt.}
\begin{itemize}
% ORIG03-ASSESSMENT-MAP-ID: diag.prompt.0001
\item \texttt{diag.prompt.0001}: norma dasar dan aritmetika vektor.
% ORIG03-ASSESSMENT-MAP-ID: diag.prompt.0002
\item \texttt{diag.prompt.0002}: rantai ketaksamaan norma.
% ORIG03-ASSESSMENT-MAP-ID: diag.prompt.0003
\item \texttt{diag.prompt.0003}: proyeksi pada hiperbidang.
% ORIG03-ASSESSMENT-MAP-ID: diag.prompt.0004
\item \texttt{diag.prompt.0004}: karakterisasi proyeksi konveks.
% ORIG03-ASSESSMENT-MAP-ID: diag.prompt.0005
\item \texttt{diag.prompt.0005}: koordinat kombinasi konveks.
% ORIG03-ASSESSMENT-MAP-ID: diag.prompt.0006
\item \texttt{diag.prompt.0006}: kestabilan kekonveksan terhadap irisan.
% ORIG03-ASSESSMENT-MAP-ID: diag.prompt.0007
\item \texttt{diag.prompt.0007}: gradien dan Hessian kuadratik.
% ORIG03-ASSESSMENT-MAP-ID: diag.prompt.0008
\item \texttt{diag.prompt.0008}: konstanta kuat-konveks dan mulus.
% ORIG03-ASSESSMENT-MAP-ID: diag.prompt.0009
\item \texttt{diag.prompt.0009}: penyangga afin orde pertama.
% ORIG03-ASSESSMENT-MAP-ID: diag.prompt.0010
\item \texttt{diag.prompt.0010}: stasioneritas dan optimalitas global.
% ORIG03-ASSESSMENT-MAP-ID: diag.prompt.0011
\item \texttt{diag.prompt.0011}: subdiferensial nilai mutlak.
% ORIG03-ASSESSMENT-MAP-ID: diag.prompt.0012
\item \texttt{diag.prompt.0012}: ambang lunak skalar.
% ORIG03-ASSESSMENT-MAP-ID: diag.prompt.0013
\item \texttt{diag.prompt.0013}: ekspektasi dan varians diskret.
% ORIG03-ASSESSMENT-MAP-ID: diag.prompt.0014
\item \texttt{diag.prompt.0014}: rata-rata sampel independen.
% ORIG03-ASSESSMENT-MAP-ID: diag.prompt.0015
\item \texttt{diag.prompt.0015}: penyelesaian rekurensi kontraktif.
% ORIG03-ASSESSMENT-MAP-ID: diag.prompt.0016
\item \texttt{diag.prompt.0016}: kompleksitas iterasi geometrik.
% ORIG03-ASSESSMENT-MAP-ID: diag.prompt.0017
\item \texttt{diag.prompt.0017}: kerucut normal dan optimalitas terkendala.
% ORIG03-ASSESSMENT-MAP-ID: diag.prompt.0018
\item \texttt{diag.prompt.0018}: bentuk titik tetap proyeksi.
% ORIG03-ASSESSMENT-MAP-ID: diag.prompt.0019
\item \texttt{diag.prompt.0019}: ortogonalitas dalam Cauchy--Schwarz.
% ORIG03-ASSESSMENT-MAP-ID: diag.prompt.0020
\item \texttt{diag.prompt.0020}: ketaksamaan Young berbobot.
\end{itemize}

\paragraph{Set soal dasar konveks: 12 prompt.}
\begin{itemize}
% ORIG03-ASSESSMENT-MAP-ID: ps01.prompt.0001
\item \texttt{ps01.prompt.0001}: maksimum fungsi afin.
% ORIG03-ASSESSMENT-MAP-ID: ps01.prompt.0002
\item \texttt{ps01.prompt.0002}: subdiferensial fungsi nilai mutlak.
% ORIG03-ASSESSMENT-MAP-ID: ps01.prompt.0003
\item \texttt{ps01.prompt.0003}: minimisasi komposit satu dimensi.
% ORIG03-ASSESSMENT-MAP-ID: ps01.prompt.0004
\item \texttt{ps01.prompt.0004}: spektrum kuadratik konveks.
% ORIG03-ASSESSMENT-MAP-ID: ps01.prompt.0005
\item \texttt{ps01.prompt.0005}: kontraksi eksak langkah gradien.
% ORIG03-ASSESSMENT-MAP-ID: ps01.prompt.0006
\item \texttt{ps01.prompt.0006}: evaluasi numerik Jensen.
% ORIG03-ASSESSMENT-MAP-ID: ps01.prompt.0007
\item \texttt{ps01.prompt.0007}: pembuktian Jensen dua titik.
% ORIG03-ASSESSMENT-MAP-ID: ps01.prompt.0008
\item \texttt{ps01.prompt.0008}: konjugat kuadratik.
% ORIG03-ASSESSMENT-MAP-ID: ps01.prompt.0009
\item \texttt{ps01.prompt.0009}: bikonjugat kuadratik.
% ORIG03-ASSESSMENT-MAP-ID: ps01.prompt.0010
\item \texttt{ps01.prompt.0010}: celah Fenchel--Young.
% ORIG03-ASSESSMENT-MAP-ID: ps01.prompt.0011
\item \texttt{ps01.prompt.0011}: proyeksi pada kotak.
% ORIG03-ASSESSMENT-MAP-ID: ps01.prompt.0012
\item \texttt{ps01.prompt.0012}: bukti pemotongan koordinat.
\end{itemize}

\paragraph{Set soal metode proksimal: 11 prompt.}
\begin{itemize}
% ORIG03-ASSESSMENT-MAP-ID: ps02.prompt.0001
\item \texttt{ps02.prompt.0001}: penurunan rumus ambang lunak.
% ORIG03-ASSESSMENT-MAP-ID: ps02.prompt.0002
\item \texttt{ps02.prompt.0002}: proks nilai mutlak berdimensi tiga.
% ORIG03-ASSESSMENT-MAP-ID: ps02.prompt.0003
\item \texttt{ps02.prompt.0003}: operator iterasi gradien proksimal.
% ORIG03-ASSESSMENT-MAP-ID: ps02.prompt.0004
\item \texttt{ps02.prompt.0004}: dua iterasi dan solusi eksak.
% ORIG03-ASSESSMENT-MAP-ID: ps02.prompt.0005
\item \texttt{ps02.prompt.0005}: lema penurunan untuk gradien Lipschitz.
% ORIG03-ASSESSMENT-MAP-ID: ps02.prompt.0006
\item \texttt{ps02.prompt.0006}: langkah kuadratik tajam.
% ORIG03-ASSESSMENT-MAP-ID: ps02.prompt.0007
\item \texttt{ps02.prompt.0007}: bentuk operator maju--mundur.
% ORIG03-ASSESSMENT-MAP-ID: ps02.prompt.0008
\item \texttt{ps02.prompt.0008}: titik tetap dan minimizer.
% ORIG03-ASSESSMENT-MAP-ID: ps02.prompt.0009
\item \texttt{ps02.prompt.0009}: kontraksi operator maju--mundur.
% ORIG03-ASSESSMENT-MAP-ID: ps02.prompt.0010
\item \texttt{ps02.prompt.0010}: resolven operator linear monoton kuat.
% ORIG03-ASSESSMENT-MAP-ID: ps02.prompt.0011
\item \texttt{ps02.prompt.0011}: laju titik proksimal.
\end{itemize}

\paragraph{Set soal dualitas dan KKT: 11 prompt.}
\begin{itemize}
% ORIG03-ASSESSMENT-MAP-ID: ps03.prompt.0001
\item \texttt{ps03.prompt.0001}: fungsi dual masalah kuadratik berekualitas.
% ORIG03-ASSESSMENT-MAP-ID: ps03.prompt.0002
\item \texttt{ps03.prompt.0002}: sistem KKT dan celah nol.
% ORIG03-ASSESSMENT-MAP-ID: ps03.prompt.0003
\item \texttt{ps03.prompt.0003}: pengali KKT kendala pertidaksamaan.
% ORIG03-ASSESSMENT-MAP-ID: ps03.prompt.0004
\item \texttt{ps03.prompt.0004}: Slater dan dual skalar.
% ORIG03-ASSESSMENT-MAP-ID: ps03.prompt.0005
\item \texttt{ps03.prompt.0005}: KKT proyeksi pada hiperbidang afin.
% ORIG03-ASSESSMENT-MAP-ID: ps03.prompt.0006
\item \texttt{ps03.prompt.0006}: nilai primal dan dual hiperbidang.
% ORIG03-ASSESSMENT-MAP-ID: ps03.prompt.0007
\item \texttt{ps03.prompt.0007}: keunikan dan kecukupan KKT.
% ORIG03-ASSESSMENT-MAP-ID: ps03.prompt.0008
\item \texttt{ps03.prompt.0008}: pengali sebagai sensitivitas nilai.
% ORIG03-ASSESSMENT-MAP-ID: ps03.prompt.0009
\item \texttt{ps03.prompt.0009}: ekspansi eksak fungsi nilai.
% ORIG03-ASSESSMENT-MAP-ID: ps03.prompt.0010
\item \texttt{ps03.prompt.0010}: KKT proyeksi pada bola.
% ORIG03-ASSESSMENT-MAP-ID: ps03.prompt.0011
\item \texttt{ps03.prompt.0011}: Slater dan kecukupan pada bola.
\end{itemize}

\subsection{Empat belas latihan terselesaikan yang sudah diterima}

Daftar berikut hanya memetakan ID dan peran kurikuler; bunyi latihan, petunjuk,
dan solusi tetap berada pada modul asalnya.
\begin{itemize}
% ORIG03-EXISTING-ASSESSMENT-ID: d90.becker.98ed693.b03.exercise.0001
\item \texttt{d90.becker.98ed693.b03.exercise.0001}: ketakbiasan dalam
  reduksi varians.
% ORIG03-EXISTING-ASSESSMENT-ID: d90.becker.98ed693.b03.exercise.0002
\item \texttt{d90.becker.98ed693.b03.exercise.0002}: identitas varians dan
  langkah kuadratik.
% ORIG03-EXISTING-ASSESSMENT-ID: d90.orig.v1.tr01.exercise.0001
\item \texttt{d90.orig.v1.tr01.exercise.0001}: operator proksimal Lasso.
% ORIG03-EXISTING-ASSESSMENT-ID: d90.orig.v1.tr01.exercise.0002
\item \texttt{d90.orig.v1.tr01.exercise.0002}: ketaksamaan satu langkah.
% ORIG03-EXISTING-ASSESSMENT-ID: d90.orig.v1.tr01.exercise.0003
\item \texttt{d90.orig.v1.tr01.exercise.0003}: sampling tanpa penggantian.
% ORIG03-EXISTING-ASSESSMENT-ID: d90.orig.v1.tr01.exercise.0004
\item \texttt{d90.orig.v1.tr01.exercise.0004}: cermin entropi.
% ORIG03-EXISTING-ASSESSMENT-ID: d90.orig.v1.tr01.exercise.0005
\item \texttt{d90.orig.v1.tr01.exercise.0005}: anggaran oracle minibatch.
% ORIG03-EXISTING-ASSESSMENT-ID: d90.orig.v1.tr01.exercise.0006
\item \texttt{d90.orig.v1.tr01.exercise.0006}: varians Prox-SAGA.
% ORIG03-EXISTING-ASSESSMENT-ID: d90.orig.v1.tr02.exercise.0001
\item \texttt{d90.orig.v1.tr02.exercise.0001}: VI dan kerucut normal.
% ORIG03-EXISTING-ASSESSMENT-ID: d90.orig.v1.tr02.exercise.0002
\item \texttt{d90.orig.v1.tr02.exercise.0002}: maksimalitas.
% ORIG03-EXISTING-ASSESSMENT-ID: d90.orig.v1.tr02.exercise.0003
\item \texttt{d90.orig.v1.tr02.exercise.0003}: resolven.
% ORIG03-EXISTING-ASSESSMENT-ID: d90.orig.v1.tr02.exercise.0004
\item \texttt{d90.orig.v1.tr02.exercise.0004}: operator skew.
% ORIG03-EXISTING-ASSESSMENT-ID: d90.orig.v1.tr02.exercise.0005
\item \texttt{d90.orig.v1.tr02.exercise.0005}: rentang langkah.
% ORIG03-EXISTING-ASSESSMENT-ID: d90.orig.v1.tr02.exercise.0006
\item \texttt{d90.orig.v1.tr02.exercise.0006}: bayangan Douglas--Rachford.
\end{itemize}

Dengan demikian, peta awal ini merujuk 54 prompt pada modul 01--04 dan 14
latihan terselesaikan yang sudah diterima tanpa mendeklarasikan ulang ID-nya.
Asesmen kumulatif lanjutan, rubrik bukti, dua laboratorium komputasi, dan
proyek kapstone kini lengkap pada modul 05--13 yang mengikuti peta ini.
